\chapter{Methodology}
% \section{Chapter overview}
% Outline the methodological structure and rationale.

\section{Research methodology}

% The following table provides an overview of the selected scientific methodology based on the
% conventional Research ‘Onion’ Model, presented by (Saunders et al., 2019).

% Table

This study adopts pragmatism as its research philosophy, as it is primarily focused on addressing a real-world problem, enhancing fundamental analysis through an Agentic AI-based approach, rather than adhering to a single philosophical paradigm. Pragmatism enables the integration of both quantitative methods and qualitative insights, allowing the research to effectively handle diverse data types. 

\section{Development methodology}

This study adopts an iterative development methodology, where the proposed Agentic AI-based system is developed in incremental stages. Each iteration involves designing and implementing a prototype, followed by systematic evaluation of its performance in supporting fundamental analysis. The insights gained from each evaluation are used to refine and enhance the subsequent version of the system. This continuous cycle of development and improvement ensures that the solution evolves progressively, leading to a more effective and reliable framework for fundamental analysis. 

\section{Project management methodology}

Agile principles are adopted for project management to support iterative development and continuous improvement, enabling the system to evolve through successive refinements until it meets the desired project objectives.

\subsection{Project Scope}

This study focuses on fundamentally analysing publicly listed companies in the Colombo Stock Exchange (CSE). However, diversified holding companies are excluded due to the complexity arising from their multi-sector operations. The analysis is primarily based on annual reports, and quarterly reports are not considered. As a result, the findings may reflect a maximum time lag of up to one year. 

\subsection{Schedule}

\begin{table}[H]
    \centering
    \begin{tabular}{|l|l|}
        \hline
        Milestone & Target date \\
        \hline
        Project Proposal Presentation  & May 7th, 2026  \\
        Project Proposal and Interim Demonstration  & May 14th, 2026  \\
        Implementation complete & June 30th, 2026  \\
        Submission  & July 30th, 2026 \\
        \hline
    \end{tabular}
    \caption{Project schedule summary}
    \label{tab:project-schedule}
\end{table}

% \subsection{Deliverables and dates}
% List key deliverables and deadlines.

% \subsection{Resources}
% \subsubsection{Hardware resources}
% List compute and infrastructure resources.

% \subsubsection{Software resources}
% List tools, frameworks, and libraries.

% \subsubsection{New technical skills}
% List technical skills acquired during the project.

% \subsection{Data requirements}
% Specify datasets, licensing, and preprocessing requirements.

% \subsection{Risks and mitigation}
% Document major project risks with likelihood, impact, and mitigation strategy.

% \begin{table}[H]
%     \centering
%     \begin{tabular}{|l|l|l|l|}
%         \hline
%         Risk & Likelihood & Impact & Mitigation \\
%         \hline
%         Data quality risk & Medium & High & Validation and cleaning pipeline \\
%         Timeline overrun & Medium & Medium & Weekly milestone tracking \\
%         \hline
%     \end{tabular}
%     \caption{Risk and mitigation table (placeholder)}
%     \label{tab:risk-mitigation}
% \end{table}

% \section{Chapter summary}
% Summarize methodological choices and transition to LESP and implementation.
